% ==============================================================================
% Appendix C: Supplementary Experiments
% "Holonomy Geometry of Cognitive States"
% ==============================================================================

\section{Supplementary Experiments}
\label{app:supplementary}

This appendix reports four supplementary experiments that verify the robustness
and interpretive structure of the main results. All experiments use the existing
180 test conversations (30 per scenario type) for Llama-2-7b unless otherwise noted.

% ------------------------------------------------------------------------------
\subsection{PCA Decomposition of Block Holonomy}
\label{app:pca}

To quantify the ``one dominant mode'' hypothesis (Section~\ref{sec:reframing}), we
perform principal component analysis on the $180 \times 5$ block holonomy matrix,
where each row is a conversation and each column is one of the five block-specific
holonomy norms $(\|F\|_{\min}, \|F\|_{\mathrm{nar}}, \|F\|_{\mathrm{soc}},
\|F\|_{\mathrm{mor}}, \|F\|_{\mathrm{pos}})$, averaged across all 32 model layers.
Columns are z-score normalized before PCA.

\paragraph{Results.}
The first principal component accounts for $99.99\%$ of total variance. The PC1
loadings are approximately uniform across all five standpoint blocks:
\begin{equation}
    \ell_{\min} = \ell_{\mathrm{nar}} = \ell_{\mathrm{soc}} = \ell_{\mathrm{mor}} =
    \ell_{\mathrm{pos}} \approx -0.447 \approx -\frac{1}{\sqrt{5}}.
\end{equation}
This confirms that PC1 captures a global coherence signal rather than any single
standpoint dimension. PC1 alone achieves $\varepsilon^2 = 0.42$ for scenario
discrimination (Kruskal--Wallallis, $p < 10^{-14}$), matching the full
five-dimensional result ($\varepsilon^2 = 0.42$--$0.68$).

Residual components PC2--PC5 collectively account for $< 0.01\%$ of variance.
Their apparent scenario discrimination after controlling for PC1 is a numerical
artifact of near-zero variance in the residual subspace; these components carry
no meaningful additional diagnostic information.

\paragraph{Interpretation.}
The five block-specific holonomy norms are linear projections of a single global
signal. The five-axis decomposition provides interpretable vocabulary (which
standpoint block receives the highest absolute load), but the primary empirical
signal is one-dimensional. This is consistent with the superposition hypothesis
in mechanistic interpretability~\cite{elhage2022superposition}: transformer
representations encode many features in high-dimensional superposition rather
than in cleanly separated subspaces.

% ------------------------------------------------------------------------------
\subsection{Baseline Ensemble Analysis}
\label{app:baseline-ensemble}

The holonomy deviation $F_{ijk}$ is defined relative to an expected transport
$U_{ik}^{\exp}$ estimated from T1 baseline conversations. This makes the measure
reference-relative rather than intrinsic. To verify that the main conclusions are
robust to baseline selection, we perform a bootstrap ensemble analysis.

\paragraph{Method.}
We randomly resample the 30 T1 conversations into five subsets of 6 conversations
each (with replacement within each subset, 20 bootstrap iterations per subset).
For each resampled baseline, we recompute the scenario discrimination (H3) and
T0 anomaly (H7) statistics.

\paragraph{Results.}
Across all five baseline subsets and 100 total bootstrap iterations:
\begin{itemize}
    \item H3~$\varepsilon^2$ ranges from $0.46$ to $0.51$ (all above the $0.30$
    robustness threshold).
    \item H7 Cohen's~$d$ ranges from $1.48$ to $3.35$ (all indicating large effect).
    \item T0 ranks first (highest holonomy) in all subsamples.
\end{itemize}

\paragraph{Conclusion.}
The scenario discrimination (H3) and T0 anomaly (H7) results are robust to
baseline selection. The T0 anomaly is not an artifact of a particular T1 baseline
choice.

% ------------------------------------------------------------------------------
\subsection{Null Grouping Controls}
\label{app:null-grouping}

The head grouping function $\gamma : H \to K$ assigns attention heads to
standpoint layers based on attention differential analysis. To verify that this
learned grouping provides non-trivial diagnostic decomposition, we compare it
against three null groupings.

\paragraph{Null grouping types.}
\begin{enumerate}
    \item \textbf{Random:} Heads randomly assigned to 5 layers, preserving the
    layer size distribution of the learned $\gamma$ (20 random instances).
    \item \textbf{Shuffled:} Random permutation of the learned $\gamma$ assignments
    (20 instances).
    \item \textbf{Layer-uniform:} Heads assigned by physical transformer layer
    position (first $1/5$ of layers $\to \psi_{\min}$, last $1/5 \to \psi_{\pos}$,
    etc.).
\end{enumerate}

\paragraph{Method.}
For each null grouping, we recompute holonomy deviation for all 180 conversations
and compute the H3~$\varepsilon^2$ scenario discrimination effect size. We then
compare the learned $\gamma$'s $\varepsilon^2$ against the null distribution using
a one-sided permutation test.

\paragraph{Status.}
This experiment requires raw activations (activations.npz) and is deferred pending
extraction. The experiment code is available in
\texttt{experiments/experiment\_b\_null\_grouping.py}. Results will be reported in
a subsequent revision.

% ------------------------------------------------------------------------------
\subsection{Competitive Baselines}
\label{app:competitive}

To verify that holonomy deviation provides unique diagnostic value beyond trivial
attention-based statistics, we compare it against three alternative measures:

\begin{enumerate}
    \item \textbf{Layer-wise attention distance:} Mean Frobenius distance between
    adjacent layer attention matrices:
    $d_{\mathrm{attn}} = \frac{1}{L-1}\sum_{l=1}^{L-1} \|A^{(l+1)} - A^{(l)}\|_F$.
    \item \textbf{Mean attention entropy:} Mean Shannon entropy of attention
    distributions across all heads and layers.
    \item \textbf{Residual stream norm change:} $\|\mathbf{h}_{x_3} -
    \mathbf{h}_{x_1}\|_2$.
\end{enumerate}

For each measure, we compute the Kruskal--Wallallis $\varepsilon^2$ for scenario
discrimination and compare against the holonomy deviation result
($\varepsilon^2 = 0.42$--$0.68$).

\paragraph{Status.}
This experiment requires raw activations and is deferred pending extraction. The
experiment code is available in
\texttt{experiments/experiment\_d\_competitive\_baselines.py}. Results will be
reported in a subsequent revision.

% ------------------------------------------------------------------------------
\subsection{Value Matrix Orthonormality}
\label{app:orthonormality}

The derivation of $V_h V_h^T = P_{\mathrm{range}(V_h)}$ (orthogonal projection)
requires $V_h^T V_h = \mathrm{Id}_{d_v}$, i.e., the columns of the value matrix
$V_h$ are orthonormal. This is not enforced by transformer training. We assess
this assumption empirically by computing, for each head~$h$, the deviation
\begin{equation}
    \epsilon_h = \|V_h^T V_h - \mathrm{Id}_{d_v}\|_F.
\end{equation}

\paragraph{Results.}
For GPT-2 (144 heads across 12 layers), the mean orthonormality deviation is
$\bar{\epsilon}_h = 82.8$ ($\sigma = 43.0$, range: $4.2$--$216.6$). Deviation
increases monotonically with layer depth: layer~0 mean $\epsilon_h = 19.4$,
layer~11 mean $\epsilon_h = 163.8$. This systematic increase suggests that
deeper layers accumulate larger deviations from orthonormality during training.

The large deviations indicate that the orthogonal projection approximation
($V_h V_h^T = P_{\mathrm{range}(V_h)}$) does not hold exactly. However, the
framework remains valid: the transport operator $U_{ij} = \sum_k
\bar{\alpha}_{ji}^{(k)} P_{W_k}$ is well-defined regardless of orthonormality,
and the block-diagonal structure (Proposition~\ref{prop:block-structure}) depends
on subspace independence ($\delta = \max_{k \neq l} \|P_{W_k} P_{W_l}\|_2$),
not on individual head orthonormality. The $\epsilon_h$ values quantify the
degree to which $V_h V_h^T$ approximates an orthogonal projector; when
$\epsilon_h$ is large, the projection $P_{W_k}$ constructed from head value
matrices is a weighted average of non-orthogonal rank-$d_v$ projections rather
than a true orthogonal projector.

Results for Llama-2-7b (1024 heads) require model weight access and will be
reported in a subsequent revision.
